\appendix
\section*{Appendix A: Formal Mathematical Foundations of the Constellation–Manifold Toolkit}

\subsection*{A.0 Altitude Structure}

We define the altitude set
\[
\mathcal{A} = 
\{A_0, A_1, \Phi_7, D_1, D_2, D_3, \Phi_8, \Phi_9, S_1, I_1\}.
\]

Each altitude $a \in \mathcal{A}$ is a purity‑restricted transformation class with:
\begin{itemize}
    \item domain type $\mathrm{Dom}(a)$,
    \item codomain type $\mathrm{Cod}(a)$,
    \item purity constraint $P(a)$,
    \item non‑generativity constraint $N(a)$,
    \item boundedness constraint $B(a)$,
    \item allowed transition relation $T(a,b)$.
\end{itemize}

The allowed transition chain is:
\[
A_0 \to A_1 \to \Phi_7 \to D_1 \to D_2 \to D_3 \to 
\Phi_8 \to \Phi_9 \to S_1 \to I_1.
\]

Forbidden transitions are defined by:
\[
\neg T(a,b) \quad \text{for all } (a,b) \notin \text{Allowed}.
\]

\subsection*{A.1 Reversible Geometry (Altitude $A_0$)}

Let $X$ be a smooth manifold.  
A reversible geometric operator is a bijection
\[
G : X \to X
\]
with inverse
\[
G^{-1} : X \to X.
\]

Purity constraint:
\[
P(A_0):\quad G \text{ does not generate new geometric structure.}
\]

Non‑generativity:
\[
N(A_0):\quad G(x) \in X \text{ only.}
\]

Boundedness:
\[
B(A_0):\quad \|G(x)\| < C.
\]

\subsection*{A.2 Conceptual Metric (Altitude $A_1$)}

A conceptual metric is a non‑negative scalar function
\[
M : X \times X \to \mathbb{R}_{\ge 0}.
\]

Purity:
\[
P(A_1):\quad M \text{ does not induce geometry.}
\]

Non‑generativity:
\[
N(A_1):\quad M \text{ produces scalars only.}
\]

\subsection*{A.3 Curvature Envelope (Altitude $\Phi_7$)}

A curvature envelope is a scalar field
\[
E : X \to \mathbb{R}.
\]

Purity:
\[
P(\Phi_7):\quad E \text{ does not modify } X.
\]

\subsection*{A.4 Diagnostics (Altitudes $D_1, D_2, D_3$)}

Collapse diagnostic:
\[
C : X \to X.
\]

Soft deformation diagnostic:
\[
S : X \to X.
\]

Safety diagnostic:
\[
\Sigma : X \to \{0,1\}.
\]

Purity:
\[
P(D_i):\quad \text{Diagnostics do not generate new structure.}
\]

Non‑generativity:
\[
N(D_i):\quad D_i(x) \in X.
\]

\subsection*{A.5 Drift Semigroup (Altitude $\Phi_8$)}

A drift semigroup is a one‑parameter family
\[
D_t : X \to X
\]
satisfying
\[
D_{t+s} = D_t \circ D_s.
\]

Boundedness:
\[
B(\Phi_8):\quad \|D_t(x)\| < C_t.
\]

\subsection*{A.6 Invariant Extraction (Altitude $\Phi_9$)}

An invariant extractor is a map
\[
I : X \to \mathbb{R}^k.
\]

Purity:
\[
P(\Phi_9):\quad I \text{ produces invariants only.}
\]

\subsection*{A.7 Steering (Altitude $S_1$)}

A steering operator is a reversible map
\[
\Theta : X \to X
\]
with inverse $\Theta^{-1}$.

Purity:
\[
P(S_1):\quad \Theta \text{ does not generate new structure.}
\]

\subsection*{A.8 Integration Layer (Altitude $I_1$)}

The integration layer produces the collision invariant:
\[
\Omega : X \to \{0,1\}.
\]

Purity:
\[
P(I_1):\quad \Omega \text{ is scalar‑valued only.}
\]

Non‑generativity:
\[
N(I_1):\quad \Omega \text{ does not modify } X.
\]

\subsection*{A.9 Altitude Transition Formalism}

Define the transition relation:
\[
T : \mathcal{A} \times \mathcal{A} \to \{0,1\}.
\]

Allowed transitions:
\[
T(A_0,A_1)=1,\;
T(A_1,\Phi_7)=1,\;
T(\Phi_7,D_1)=1,\;
T(D_1,D_2)=1,\;
T(D_2,D_3)=1,
\]
\[
T(D_3,\Phi_8)=1,\;
T(\Phi_8,\Phi_9)=1,\;
T(\Phi_9,S_1)=1,\;
T(S_1,I_1)=1.
\]

Forbidden transitions:
\[
T(a,b)=0 \quad \text{otherwise}.
\]

\subsection*{A.10 Firewall Constraint}

The Series V firewall enforces:
\[
\text{Firewall}(a,b,f) =
\begin{cases}
f(), & T(a,b)=1 \text{ and } P(b) \land N(b) \land B(b), \\
\text{Violation}, & \text{otherwise}.
\end{cases}
\]

This binds Lean proofs to Julia runtime execution.

\subsection*{A.11 Non‑Generativity Theorem}

For all altitudes $a$:
\[
N(a) \implies \forall x \in X,\; f_a(x) \in \mathrm{Cod}(a) \subseteq X.
\]

\subsection*{A.12 Purity Theorem}

For all altitudes $a$:
\[
P(a) \implies f_a \text{ does not induce new structure.}
\]

\subsection*{A.13 Boundedness Theorem}

For all altitudes $a$:
\[
B(a) \implies \|f_a(x)\| < C_a.
\]

\subsection*{A.14 Collision Invariant Correctness}

The integration layer satisfies:
\[
\Omega(x) = 1 \iff \text{system is collision‑safe}.
\]

\[
\Omega(x) = 0 \iff \text{violation detected}.
\]

This completes the formal mathematical foundation of the altitude architecture.
